\documentclass[11pt]{article}
\usepackage[margin=0.98in]{geometry}
\usepackage{fontspec,amsmath,amssymb,amsthm}
\usepackage[hidelinks,pdftitle={A Sharp Comparison of a Truncated Quasinorm and an Inverse Peetre Functional},pdfauthor={ProofRush Research Team}]{hyperref}
% Small OFL-licensed fallback included with this source.
\IfFileExists{fonts/ProofRushDirectorJP.otf}
  {\newfontfamily\namefont[Path=fonts/]{ProofRushDirectorJP.otf}}
  {\IfFontExistsTF{Noto Sans CJK JP}
    {\newfontfamily\namefont{Noto Sans CJK JP}}
    {\IfFontExistsTF{Hiragino Sans}
      {\newfontfamily\namefont{Hiragino Sans}}
      {\newfontfamily\namefont{HaranoAjiMincho-Regular.otf}}}}
\newtheorem{theorem}{Theorem}
\newcommand{\director}{{\namefont 大束良明} (Yoshiaki Otsuka)}
\newcommand{\F}{F_p}
\newcommand{\D}{D_p}
\title{A Sharp Comparison of a Truncated Quasinorm\\ and an Inverse Peetre Functional}
\author{ProofRush Research Team\\\small Research directed by \director}
\date{27 November 2026 JST (26 November 2026 UTC)\\\small Preprint, version 1}
\begin{document}
\maketitle

\begin{abstract}
We prove a direct comparison between the truncated \(\ell_{2/3}\)
quasinorm in the Valiant--Valiant identity-testing bounds and the inverse
\((\ell_1,\ell_2)\) Peetre functional in the Blais--Canonne--Gur bounds.
This answers the inequalities in Clément Canonne's 2019 MathOverflow
question and gives a direct answer to his COLT 2024 Open Question~2.
The lower comparison holds at the same radius. The upper comparison
uses a smaller truncation radius and has an optimal coefficient \(1/4\)
in its radius-gap term. The upper comparison cannot hold at the same
radius, even for probability vectors of arbitrarily small
\(\ell_2\)-norm. The proof uses deterministic inequalities on finite
or countably infinite probability laws.
\end{abstract}

\section{The question and the bounds}
For a probability sequence \(p_1\ge p_2\ge\cdots\ge0\), let
\[
 N_e=\min\Bigl\{N\ge0:\sum_{i=1}^N p_i\ge1-e\Bigr\},
 \qquad
 \F(e)=\left(\sum_{i=2}^{N_e}p_i^{2/3}\right)^{3/2}
 \quad (0<e\le1).
\]
The retained prefix includes the entire atom at the cutoff. The sum is
zero when \(N_e\le1\). In the notation of Canonne's MathOverflow question,
\(\F(e)=\Phi(e,p)\) \cite{mo}. Define
\[
 K_p(t)=\inf_{\substack{a\in\ell_1,\,b\in\ell_2\\a+b=p}}
                 \bigl(\|a\|_1+t\|b\|_2\bigr),\qquad
 \D(e)=\inf\{t\ge0:K_p(t)\ge1-e\}.
\]
The infimum over decompositions is over real signed sequences. In
\cite{mo}, \(\D(e)=\Psi(e,p)\) for \(0<e<1\). Canonne's COLT 2024
formulation uses \(K_p^{-1}(1-2e)\) for its BCG quantity \cite{colt}.

The testing bounds of Valiant and Valiant \cite{vv} and Blais, Canonne,
and Gur \cite{bcg} yield the following comparison by combining
different sample-complexity estimates. Canonne asked for a direct
argument. His equation~(4) is, for all \(p\) and \(0<e<1/9\),
\begin{equation}\label{moeq}
 1\vee\Phi(9e,p)\ \lesssim\ 1\vee\Psi(e,p)\
   \lesssim\ 1\vee \frac{\Phi(e/32,p)}{e}.
\end{equation}
Here \(\lesssim\) allows an absolute constant factor.

\begin{theorem}\label{main}
For every probability sequence \(p\) and \(0<s<e<1\),
\begin{equation}\label{ineq}
 \F(e)\le\D(e)\le
 \sqrt{1+\left(\frac{\F(s)}{4(e-s)}\right)^2}.
\end{equation}
The left coefficient \(1\) and the coefficient \(1/4\) in the
right-hand radius-gap term are optimal universal constants.
For each fixed \(0<e<1\), there exist laws for which \(\F(e)\) is
bounded and \(\D(e)\) is unbounded. These sharpness and obstruction
claims hold even on finite laws under any prescribed positive upper
bound on \(\|p\|_2\).
\end{theorem}

In \eqref{moeq}, the left constant can be taken to be \(1\) by monotonicity of
\(\F\). Choosing \(s=e/32\) gives
\(\D(e)\le1+(8/31)\F(e/32)/e\), so the right constant may be taken to
be \(39/31\). If \(\|p\|_2\le1/2\), the bounds
\(\D(e)\ge(1-e)/\|p\|_2\) and
\(\F(e/32)\ge1-e/32-p_1\) make both right-hand baseline terms
inactive in this range.

\section{Proof of the comparisons}
\paragraph{Inverse representation.}
For \(0<e<1\), \(\D(e)\) is the minimum \(\ell_2\)-norm of a vector
\(0\le z_i\le1\) with \(\sum_i p_i z_i\ge1-e\).
Indeed, every such \(z\in\ell_2\) satisfies
\(K_p(\|z\|_2)\ge\sum_i p_i z_i\) by Hölder's inequality applied to
each signed decomposition. To see that equality of the inverse and
the minimum holds, choose \(h>0\) so that \(z_i=\min(1,hp_i)\) satisfies
\(\sum_i p_i z_i=1-e\). Such \(h\) exists by continuity and dominated
convergence as the captured mass increases from zero to one.
Set \(b=z/h\), \(a=p-b\). Since \(a_i z_i=a_i\), the target is
\(\|a\|_1+\|z\|_2^2/h\). At every \(t<\|z\|_2\) this decomposition
has cost below the target. This gives the inverse formula with
attainment on a countable alphabet. It is a special case of the
standard dual formula for \(K_p\) \cite{ms}.

The function \(K_p\) is concave and nondecreasing by definition. It
tends to one, since the vectors \(\min(1,hp_i)\) have captured mass
tending to one and \(K_p\le1\). Concavity implies strict increase
where the value is below one, so no inverse convention affects
\(\D(e)\) for \(0<e<1\). At \(e=1\) one has \(\F(1)=\D(1)=0\)
with \(K_p(0)=0\). For completeness, if \(p\) has \(m\) positive
coordinates, the top level forms the plateau \(t\ge\sqrt m\):
the indicator of the support attains mass one there. For any
integer \(r>t^2\) positive coordinates, placing a common positive
mass \(c\) from each into \(b\) gives cost \(1-cr+tc\sqrt r<1\)
when \(c\) is at most their minimum. This excludes earlier
plateaus and excludes attainment of level one at finite \(t\)
for infinite support.

\paragraph{Lower comparison.}
If \(N=N_e=1\), the claim is immediate. Otherwise \(p_N>0\). For
any feasible \(z\), minimality of \(N\) gives
\(\sum_{i\ge N}p_i z_i\ge\sum_{i<N}p_i(1-z_i)\).
On \(i\ge N\) we have \(p_i^{2/3}\ge p_i/p_N^{1/3}\), and on
\(i<N\) the inequality reverses. Thus
\[
 \sum_{i=2}^N p_i^{2/3}
 \le\sum_{i<N}p_i^{2/3}
 \le\sum_i p_i^{2/3}z_i
 \le
 \left(\sum_i p_i\right)^{2/3}\left(\sum_i z_i^3\right)^{1/3}
 \le\left(\sum_i z_i^2\right)^{1/3}.
\]
The penultimate step is Hölder's inequality. Raise to the power
\(3/2\) and minimize over \(z\).

\paragraph{Upper comparison.}
Protect the largest atom by choosing \(z_1=1\) and put
\(S=\sum_{i=2}^{N_s}p_i^{2/3}\). If \(S=0\), this coordinate alone
suffices. Otherwise set \(\lambda=4(e-s)/S\) and choose
\(z_i=\min(1,p_i^{1/3}/\lambda)\) for \(2\le i\le N_s\) and zero
after the cutoff. For \(x\ge0\),
\(x(1-x/\lambda)_+\le\lambda/4\), where \(u_+=\max(u,0)\).
Writing \(x=p_i^{1/3}\) for each retained atom gives
\[
 \sum_{i=2}^{N_s}p_i(1-z_i)\le S\lambda/4=e-s,\qquad
 \sum_{i=2}^{N_s}z_i^2\le S/\lambda^2.
\]
The discarded mass is at most \(s\), so the witness is feasible
and has squared norm at most \(1+\F(s)^2/[16(e-s)^2]\).

\section{Sharpness and the radius change}
For the uniform law on \(m\) labels, symmetry or Cauchy--Schwarz
gives \(\D(e)=(1-e)\sqrt m\), and
\(\F(e)=(\lceil(1-e)m\rceil-1)^{3/2}/m\).
First \(m\to\infty\) and then \(e\downarrow0\) makes
\(\F(e)/\D(e)\) approach \(1\).

For sharpness of the upper coefficient fix \(0<s<e<1\) with
\(2e-s<1\), and set \(g=e-s\) and \(H=1-s-2g\).
Consider \(L\) equal head atoms of total mass \(H\), followed by
\(m\) equal atoms of total mass \(2g\) and \(m^2\) equal atoms of
total mass \(s\). Here \(L\) is fixed and \(m\) is large, so the
blocks are in decreasing order. The minimal prefix at radius \(s\)
ends with the second block, and \(\F(s)/(2g\sqrt m)\to1\).
Every feasible witness satisfies
\[
 H+g
 \le H+\left(2g/\sqrt m+s/m\right)\|z\|_2 .
\]
This follows by bounding the captured mass of the first block by
\(H\) and applying Cauchy--Schwarz to each of the other blocks.
It yields \(\liminf_{m\to\infty}g\,\D(e)/\F(s)\ge1/4\).
The upper bound proves equality in the limit. Taking \(L\) large
and then \(m\) large gives any prescribed positive \(\ell_2\) ceiling.

For the same-radius obstruction, fix \(0<e<1\), set
\(\gamma=m^{-1/3}\), and discard small \(m\) for which the
construction is unordered or has a nonpositive head. The law
has \(L\) equal head atoms of total mass \(1-e-\gamma\) and \(m\)
equal tail atoms of total mass \(t=e+\gamma\).
The cutoff keeps \(\lceil m\gamma/t\rceil\) tail atoms. Their
contribution to \(\F(e)^{2/3}\) is bounded by
\((m\gamma/t+1)(t/m)^{2/3}=O_e(1)\). The head contribution is also
bounded at fixed \(L\). A feasible witness must capture at least
\(\gamma\) from the tail, and Cauchy--Schwarz gives
\(\D(e)\ge\gamma\sqrt m/t\to\infty\).
Choosing \(L\) sufficiently large before sending \(m\) to infinity
again meets any prescribed positive \(\ell_2\) ceiling.

\section{Relation to prior work and provenance}
The dual description of \(K_p\) and its connection to Rademacher
sums are classical \cite{ms}. Holmstedt's estimates and the
partition \(Q\)-norm comparison give useful approximations to
\(K_p\) itself \cite{holm,bcg}. A multiplicative approximation
to \(K_p\) can shift the inverse level near one. The present
argument treats that level directly. The point-mass example in
Pinelis's answer to \cite{mo} explains the need for a baseline in a
universal upper comparison. A targeted literature search did not
find a prior direct proof of \eqref{moeq} or \eqref{ineq}.
This search gives limited evidence on priority. External peer
review is pending. Section~5 of \cite{colt} names an extra-large
wombat plush toy for Question~2. The AUD~200 offer concerns
Question~1 on minimax sample complexity.

\paragraph{Contributions and computation.}
\director\ directed the research and set its criteria.
OpenAI Codex agents served as the ProofRush mathematical team.
GPT-6.2-Astra at maximum reasoning effort coordinated and developed
the analytic proof. GPT-6.2-Sol at medium effort independently
reviewed the proof, reproduced the bounds, checked examples, and
drafted this manuscript. GPT-6.1-Terra at medium effort audited the
prior literature and formalized lower-bound steps. GPT-6.2-Sol at
high effort worked on formalizing the upper witness and reconciling
the source radii. Exact arithmetic checks on an Apple M3 Pro
covered 1,000 finite examples, 20 cutoff-boundary cases, and one
countable example with divergent untruncated
\(\ell_{2/3}\) quasinorm. The checks supplement the analytic proof.
No additional paid API was used.

\paragraph{Formal verification and reproduction.}
The \href{https://github.com/yoshi0703/ProofRush/tree/canonne-functional-comparison-v0.1.0/studies/canonne-functional-comparison}{companion archive}
contains the Lean~4.36 proofs, pinned Mathlib revision, source-matching
audit, and replay commands. Lean and the separate Rust kernel nanoda
verified the inverse level for countable laws, both inequalities in
\eqref{ineq}, the original radii in \eqref{moeq}, optimality of both
universal coefficients, and the same-radius obstruction at every fixed
interior radius. The split-head refinements imposing arbitrarily small
\(\ell_2\)-norm are proved above.

\begin{thebibliography}{9}\small
\bibitem{mo} C. Canonne, \emph{Relation between two different
functionals: \(\lVert p^{-\max}_{-\varepsilon}\rVert\) and
\(\kappa_p^{-1}(1-\varepsilon)\)}, MathOverflow question 332062
(20 May 2019), \href{https://mathoverflow.net/questions/332062/relation-between-two-different-functionals-lvert-p-max-varepsilon-rve}{question and discussion}.
\bibitem{colt} C. Canonne, \emph{Open Problem: Tight Characterization of
Instance-Optimal Identity Testing}, COLT 2024, PMLR 247, 5312--5316.
\href{https://proceedings.mlr.press/v247/canonne24a.html}{Proceedings}.
\bibitem{vv} G. Valiant and P. Valiant, \emph{An Automatic Inequality
Prover and Instance Optimal Identity Testing}, SIAM J. Comput.
46(1) (2017), 429--455. \href{https://doi.org/10.1137/151002526}{DOI}.
\bibitem{bcg} E. Blais, C. L. Canonne, and T. Gur, \emph{Distribution
Testing Lower Bounds via Reductions from Communication Complexity},
ACM Trans. Comput. Theory 11(2) (2019), article 6, pp. 6:1--6:37.
\href{https://doi.org/10.1145/3305270}{Journal}
and \href{https://doi.org/10.4230/LIPIcs.CCC.2017.28}{CCC 2017, LIPIcs 79, pp. 28:1--28:40}.
\bibitem{ms} S. J. Montgomery-Smith, \emph{The distribution of Rademacher
sums}, Proc. Amer. Math. Soc. 109(2) (1990), 517--522.
\href{https://doi.org/10.1090/S0002-9939-1990-1013975-0}{DOI}.
\bibitem{holm} T. Holmstedt, \emph{Interpolation of quasi-normed
spaces}, Math. Scand. 26 (1970), 177--199.
\href{https://doi.org/10.7146/math.scand.a-10976}{DOI}.
\end{thebibliography}
\end{document}
